\section{The process}
\label{bdc:sec:process}

The birth--death--catastrophe process was introduced in \ChM{} as
\mextref{m:def:bdc}{Definition 2.4}. We restate it here for the single-cell
reading.

\begin{definition}[Birth--death--catastrophe process]
\label{bdc:def:bdc}
Let $X_t$ take values in $\{0,1,2,\ldots\}\cup\{H\}$ and record the
intracellular load while the cell remains intact. Each replicative unit
independently gives birth at rate $\lambda$, dies at rate $\mu$, and
triggers rupture at rate $\delta$. For $n\ge1$, the non-zero transition
rates are
\begin{equation}
\label{bdc:eq:rates}
q_{n,n+1}=\lambda n,
\qquad
q_{n,n-1}=\mu n,
\qquad
q_{n,H}=\delta n,
\end{equation}
and both $0$ and $H$ are absorbing.
\end{definition}

The units counted by $X_t$ are infectious particles inside a host cell, and
rupture releases the entire load at once. Because the total catastrophe rate
in state $n$ is $\delta n$, a crowded cell is more likely to burst than a
sparsely occupied one. This population dependence is the small addition that
changes the character of the process. The same chain is drawn beside the
linear birth--death process in \mextref{m:fig:ratediagram}{Figure 2.6 of
\ChM}; \cref{bdc:fig:ratediagram} repeats it here so that the two absorbing
mechanisms remain in view.

\begin{figure}[!htbp]
\centering
\resizebox{0.92\textwidth}{!}{%
\begin{tikzpicture}[
  >={Stealth[length=1.9mm]},
  st/.style={circle,draw,thick,minimum size=7.5mm,inner sep=0pt,font=\small},
  ab/.style={circle,draw,thick,double,double distance=0.7pt,
             minimum size=7.5mm,inner sep=0pt,font=\small},
  lb/.style={font=\footnotesize,inner sep=1.6pt},
  rowlab/.style={font=\small\itshape,anchor=west},
  every path/.style={thick}
]
\node[rowlab] at (-1.9,1.15) {birth--death};
\node[ab] (a0) at (0,0) {$0$};
\node[st] (a1) at (2.1,0) {$1$};
\node[st] (a2) at (4.2,0) {$2$};
\node[st] (a3) at (6.3,0) {$3$};
\node (a4) at (7.9,0) {$\cdots$};
\draw[->] (a1) to[bend left=38] node[above,lb]{$\lambda$} (a2);
\draw[->] (a2) to[bend left=38] node[above,lb]{$2\lambda$} (a3);
\draw[->] (a3) to[bend left=38] node[above,lb]{$3\lambda$} (a4);
\draw[->] (a1) to[bend left=38] node[below,lb]{$\mu$} (a0);
\draw[->] (a2) to[bend left=38] node[below,lb]{$2\mu$} (a1);
\draw[->] (a3) to[bend left=38] node[below,lb]{$3\mu$} (a2);

\begin{scope}[yshift=-3.5cm]
\node[rowlab] at (-1.9,1.15) {birth--death--catastrophe};
\node[ab] (b0) at (0,0) {$0$};
\node[st] (b1) at (2.1,0) {$1$};
\node[st] (b2) at (4.2,0) {$2$};
\node[st] (b3) at (6.3,0) {$3$};
\node (b4) at (7.9,0) {$\cdots$};
\draw[->] (b1) to[bend left=38] node[above,lb]{$\lambda$} (b2);
\draw[->] (b2) to[bend left=38] node[above,lb]{$2\lambda$} (b3);
\draw[->] (b3) to[bend left=38] node[above,lb]{$3\lambda$} (b4);
\draw[->] (b1) to[bend left=38] node[below,lb]{$\mu$} (b0);
\draw[->] (b2) to[bend left=38] node[below,lb]{$2\mu$} (b1);
\draw[->] (b3) to[bend left=38] node[below,lb]{$3\mu$} (b2);
\node[ab] (bh) at (4.2,-2.35) {$H$};
\draw[->,densely dashed] (b1) to[bend right=12]
  node[pos=0.68,left=1.5pt,lb]{$\delta$} (bh);
\draw[->,densely dashed] (b2) --
  node[pos=0.68,right=1.5pt,lb]{$2\delta$} (bh);
\draw[->,densely dashed] (b3) to[bend left=12]
  node[pos=0.68,right=1.5pt,lb]{$3\delta$} (bh);
\end{scope}
\end{tikzpicture}}
\caption{Linear birth--death dynamics (top) and the
birth--death--catastrophe process used throughout this chapter (bottom).
Double circles are absorbing. Ordinary extinction enters $0$ only from
state $1$; catastrophe jumps from any positive count to the distinct state
$H$, at total rate $\delta n$.}
\label{bdc:fig:ratediagram}
\end{figure}

\subsection{Sample paths}
\label{bdc:sec:paths}

Unless a caption states otherwise, figures use the working point
$(\lambda,\mu,\delta)=(1,0.2,0.05)$, with $X_0=1$.

\begin{figure}[!htbp]
\centering
\includegraphics[width=0.78\textwidth]{F4_paths_XW.pdf}
\caption{Ten exact realisations of the load $X_t$ (top) and the released
count $W_t$ (bottom), generated by the Gillespie algorithm
\cite{Gillespie1977}. The step of $X_t$ to the plotting baseline marks
entry into $H$, not the count state $0$. Each $W_t$ path remains at zero
until rupture, then jumps to the load present at that instant.
$\lambda=1$, $\mu=0.2$, $\delta=0.05$, $X_0=1$; seed 12.}
\label{bdc:fig:gillespie}
\end{figure}

\begin{figure}[!htbp]
\centering
\includegraphics[width=0.78\textwidth]{F4_fate_coloured_paths.pdf}
\caption{Twelve exact realisations of $X_t$, coloured by terminal fate.
Grey paths ending in crosses reach internal extinction; vermillion paths
ending in filled circles enter $H$. The bold curve is the exact mean
$J(t)$. $\lambda=1$, $\mu=0.2$, $\delta=0.05$, $X_0=1$; seed 154.}
\label{bdc:fig:fatepaths}
\end{figure}

\Cref{bdc:fig:gillespie} shows a handful of exact realisations of the pair
$(X_t,W_t)$. Some paths decline to $0$; others grow for a time and end
abruptly at $H$. The corresponding release paths contain only one possible
movement: a jump from zero to the load carried at rupture.

Colouring the paths by their eventual fate, as in \cref{bdc:fig:fatepaths},
makes visible what a single undifferentiated bundle conceals. The exact mean
$J(t)=\ex{X_t}$ passes through the ensemble, but it belongs to neither class
of path.

\begin{remark}[The mean load describes no typical path]
\label{bdc:rem:Jmixes}
After a realisation reaches $0$ or $H$, its numerical load is zero. The
late-time mean therefore averages substantial loads from a shrinking set of
productive paths against zeros from every path that has finished. It is neither
the load of a typical cell nor the load of a typical surviving cell. The latter
is the conditional mean $J/\Ifix$, considered in
\cref{bdc:sec:plateau}.
\end{remark}

\FloatBarrier

\subsection{The released count and the joint process}
\label{bdc:sec:released}

Rupture creates a second count, the number of particles released into the
extracellular medium. It remains zero while the cell is intact and, in the
absence of any imposed extracellular dynamics, remains fixed after release. Let
\[
\tau=\inf\{t:X_t=H\}
\]
be the catastrophe time. The released count is
\begin{equation}
\label{bdc:eq:release}
W_t=
\begin{cases}
0,&t<\tau,\\
X_{\tau-},&t\ge\tau.
\end{cases}
\end{equation}
The burst size is the load present an instant before rupture,
$\Kb:=X_{\tau-}$, as recorded in
\mextref{m:eq:bdcrelease}{equation 2.20 of \ChM}. The left limit matters: at
$\tau$ the internal state has become $H$, while the released load is the
integer present an instant earlier.

The joint transition list makes the relation between the two counts explicit.
Writing $\Delta X_t=X_{t+\Delta t}-X_t$ and similarly for $W_t$, and conditioning
on $X_t=n\ge1$,
\begin{align}
\label{bdc:eq:joint}
\pr{\Delta X_t = 1,\ \Delta W_t = 0 \mid X_t=n} &=
\lambda n\,\Delta t+o(\Delta t)
&&\text{(birth)}\notag\\
\pr{\Delta X_t = -1,\ \Delta W_t = 0 \mid X_t=n} &=
\mu n\,\Delta t+o(\Delta t)
&&\text{(death)}\notag\\
\pr{X_{t+\Delta t} = H,\ W_{t+\Delta t} = n \mid X_t=n} &=
\delta n\,\Delta t+o(\Delta t)
&&\text{(catastrophe)}\\
\pr{\Delta X_t = 0,\ \Delta W_t = 0 \mid X_t=n} &=
1-(\lambda+\mu+\delta)n\,\Delta t+o(\Delta t)
&&\text{(no event)}\notag .
\end{align}
This pair $(X_t,W_t)$ will be followed through its means, variances and
covariance. Most of the work lies in the marginal moments; the joint structure,
as we shall see, is unusually simple.

\subsection{Two conventions for rupture}
\label{bdc:sec:conventions}

There is another reasonable way to represent rupture: send the process to $0$
rather than to a separate state $H$. After all, once a cell has ruptured it
contains no replicative units, although the cell itself has also ceased to
exist. We call the transition to $0$ \emph{killing} and the transition to $H$
\emph{catastrophe}. \Cref{bdc:fig:killing} shows the distinction.

\begin{figure}[H]
\centering
\includegraphics[width=0.96\textwidth]{F4_killing_vs_catastrophe.pdf}
\caption{Two conventions for rupture. Under killing (left), rupture from
state $n$ ends at $0$ and is indistinguishable from internal extinction.
Under catastrophe (right), used here, the same transition ends at the
separate state $H$. Birth and death rates remain $\lambda n$ and $\mu n$.}
\label{bdc:fig:killing}
\end{figure}

The choice does not alter the observed burst times or burst sizes. It alters what
can be distinguished after absorption. Under killing, internal extinction and
rupture arrive at the same state; their probabilities cannot be recovered
separately from the state occupancy. With $H$ kept apart, we may distinguish the
internal-extinction probability $D$, the rupture probability $1-I$, and the
productive probability $\Ifix=I-D$. That distinction is the reason for using
$H$ throughout. \fwd{singlecell}{Chapter~5} carries both conventions through
the generating-function treatment.

\subsection{Notation}
\label{bdc:sec:notation}

Although $H$ is a label rather than a number, after rupture the intracellular
\emph{load} is zero. Whenever $X_t$ appears arithmetically---in a sum, product
or moment---we use the numerical-load convention
\[
X_t=n\mapsto n,\qquad X_t=H\mapsto0.
\]
Thus $H$ remains distinct from $0$ as a state, allowing the two absorbing
mechanisms to be distinguished, but contributes zero to
$\ex{X_t^p}$ and to products such as $X_tW_t$.

One possible collision is best removed now: $\Kb$ denotes the burst size,
whereas $K(t)=\ex{X_t^2}$ will denote the second moment of the load. They are
unrelated. The full list of symbols appears in \cref{bdc:tab:notation}. Unless
explicitly stated otherwise, the process begins with one unit, $X_0=1$.

\FloatBarrier
